% P11 paper module: R33 explicit P02-to-P11 archimedean Gamma symbol bridge.
% Constructive identification of g_infty with the P02 digamma symbol
% and displayed affine relation m_Gamma = c_0 + c_1 q_Gamma with c_1 != 0.

\subsection{Explicit P02-to-P11 Gamma symbol bridge}
\label{subsec:o3af-symbol-bridge}

R31 records (Remark~\ref{rem:o3ad-symbol-interface}) that the anti-locality
theorem for the exact P02 Gamma symbol transfers to the concrete P11 operator
\(C_{\Gamma,S}\) only after the affine identification
\(m_\Gamma=c_0+c_1q_\Gamma\), \(c_1\ne0\), is made explicit.  This subsection
performs that identification constructively.  It records neither Object~X nor
polar-gauge nor terminal-transport nor RH consequences; it is a manuscript-
level interface closure.

\subsubsection*{The frozen P02 diagonal symbol}

P02 fixes the archimedean logarithmic derivative
\[
 \gamma_\infty(t)
 :=-\tfrac12\log\pi
 +\tfrac12\psi\!\left(\tfrac14+\tfrac{it}{2}\right)
\]
and the even diagonal symbol
\[
 q_\Gamma(\xi)
 :=2\operatorname{Re}\gamma_\infty(\xi)
 =-\log\pi
 +\operatorname{Re}\psi\!\left(\tfrac14+\tfrac{i\xi}{2}\right).
\tag{O3AF.1}\label{eq:o3af-qgamma-def}
\]
The P02 Fourier convention
\(\widehat f(\xi)=\int f(u)e^{i\xi u}du\),
\(f(u)=(2\pi)^{-1}\!\int\widehat f(\xi)e^{-i\xi u}d\xi\)
is used throughout.

\begin{lemma}[Convergent Digamma difference]
\label{lem:o3af-digamma-difference}
For every \(\xi\in\mathbb R\),
\[
 \boxed{
 q_\Gamma(\xi)-q_\Gamma(0)
 =\sum_{n=0}^\infty
 \left[\frac1{n+\tfrac14}-\frac{n+\tfrac14}{(n+\tfrac14)^2+\tfrac{\xi^2}{4}}\right].
 }
\tag{O3AF.2}\label{eq:o3af-digamma-difference}
\]
The summands are \(O(n^{-3})\) at fixed \(\xi\), so the series converges
absolutely.  The right-hand side is nonnegative and vanishes iff \(\xi=0\).
\end{lemma}

\begin{proof}
The absolutely convergent representation of \(\psi\) on
\(\operatorname{Re}z>0\) reads
\[
 \psi(z)=-\gamma_{\rm E}+\sum_{n=0}^\infty
 \left(\frac1{n+1}-\frac1{n+z}\right).
\]
Set \(z=\tfrac14+\tfrac{i\xi}{2}\) and take real parts:
\[
 \operatorname{Re}\psi(z)
 =-\gamma_{\rm E}+\sum_{n=0}^\infty
 \left(\frac1{n+1}-\frac{n+\tfrac14}{(n+\tfrac14)^2+\tfrac{\xi^2}{4}}\right).
\]
Subtracting the same formula at \(\xi=0\) removes both \(-\gamma_{\rm E}\)
and the \(1/(n+1)\)-terms, leaving \eqref{eq:o3af-digamma-difference}.
Nonnegativity is termwise:
\(1/(n+\tfrac14)\ge(n+\tfrac14)/((n+\tfrac14)^2+\tfrac{\xi^2}{4})\)
with equality iff \(\xi=0\).
\end{proof}

\subsubsection*{Explicit construction of the P11 log symbol \(g_\infty\)}

The current P11 (§2, equation~\eqref{eq:gamma-symbol}) records only
\[
 m_\Gamma(\xi)=1+g_\infty(\xi),
 \qquad m_\Gamma(\xi)\asymp\log(2+|\xi|),
\]
without an explicit formula for \(g_\infty\).  We now fix that formula.

\begin{definition}[Explicit P11 log symbol]
\label{def:o3af-ginfty}
\[
 \boxed{
 g_\infty(\xi)
 :=q_\Gamma(\xi)-q_\Gamma(0),
 }
\tag{O3AF.3}\label{eq:o3af-ginfty}
\]
with \(q_\Gamma\) from~\eqref{eq:o3af-qgamma-def}.
\end{definition}

By Lemma~\ref{lem:o3af-digamma-difference}, \(g_\infty(\xi)\ge0\) and
\(g_\infty(0)=0\).  This is the ``explicit positive series for \(g_\infty\)''
already referenced in Subsection~\ref{subsec:o3ac-riesz-support}
(cf.\ \(\phi_S\) construction) and in the direct-terminal-bridge module.

\begin{proposition}[Symbol properties under Definition~\ref{def:o3af-ginfty}]
\label{prop:o3af-symbol-props}
The multiplier \(m_\Gamma=1+g_\infty\) satisfies:
\begin{enumerate}
\item[(i)] \(m_\Gamma:\mathbb R\to[1,\infty)\), even, real-analytic;
\item[(ii)] \(m_\Gamma(\xi)=\log|\xi|+\bigl(1-\log\pi+\psi(\tfrac14)-\log 2\bigr)
   +o(1)\) as \(|\xi|\to\infty\);
   in particular \(m_\Gamma(\xi)\asymp\log(2+|\xi|)\).
\item[(iii)] The dilation inequality \(m_\Gamma(\lambda\xi)\le\lambda^2\,m_\Gamma(\xi)\)
   for \(\lambda\ge1\) required in the direct-terminal-bridge (equation
   DT.25) holds.
\end{enumerate}
\end{proposition}

\begin{proof}
(i) By Lemma~\ref{lem:o3af-digamma-difference},
\(g_\infty\ge0\) with \(g_\infty(0)=0\), so \(m_\Gamma\ge1\).  Evenness follows
from \(\overline{\gamma_\infty(t)}=\gamma_\infty(-t)\).  Real-analyticity is
inherited from \(\operatorname{Re}\psi(\tfrac14+i\xi/2)\).

(ii) The Digamma asymptotic
\(\psi(z)=\log z-\tfrac{1}{2z}+O(1/z^2)\) as \(|z|\to\infty\) in any sector
avoiding the negative real axis gives, for \(z=\tfrac14+i\xi/2\) with real
\(\xi\to\pm\infty\),
\[
 \operatorname{Re}\psi(z)=\log|z|+O(1/|\xi|^2)=\log(|\xi|/2)+O(1/|\xi|^2).
\]
Substituting into~\eqref{eq:o3af-qgamma-def} and using
\(q_\Gamma(0)=\psi(\tfrac14)-\log\pi\),
\begin{align*}
 m_\Gamma(\xi)
 &=1+q_\Gamma(\xi)-q_\Gamma(0)\\
 &=1+\bigl(\log(|\xi|/2)-\log\pi+o(1)\bigr)-\bigl(\psi(\tfrac14)-\log\pi\bigr)\\
 &=\log|\xi|+\bigl(1-\log 2-\psi(\tfrac14)\bigr)+o(1).
\end{align*}
Both sides are bounded between constant multiples of \(\log(2+|\xi|)\) for
\(|\xi|\ge1\); at \(|\xi|\le1\) \(m_\Gamma\) is continuous, real-analytic
and bounded away from \(0\) and \(\infty\).  Together this proves
\(m_\Gamma\asymp\log(2+|\xi|)\).

(iii) Set \(y_n:=n+\tfrac14\) and \(x:=|\xi|/2\).  Then
\[
 s_n(\xi):=\frac1{y_n}-\frac{y_n}{y_n^2+x^2}
 =\frac{x^2}{y_n(y_n^2+x^2)},
 \qquad
 g_\infty(\xi)=\sum_{n\ge0}s_n(\xi).
\]
For \(\lambda\ge1\),
\[
 \frac{s_n(\lambda\xi)}{s_n(\xi)}
 =\lambda^2\cdot\frac{y_n^2+x^2}{y_n^2+\lambda^2x^2}
 \le\lambda^2,
\]
since \(y_n^2+\lambda^2x^2\ge y_n^2+x^2\).  Summing over \(n\) gives
\(g_\infty(\lambda\xi)\le\lambda^2\,g_\infty(\xi)\).
Because \(1\le\lambda^2\) as well,
\[
 m_\Gamma(\lambda\xi)
 =1+g_\infty(\lambda\xi)
 \le\lambda^2+\lambda^2g_\infty(\xi)
 =\lambda^2\,m_\Gamma(\xi).
\]
\end{proof}

\subsubsection*{The affine bridge}

\begin{corollary}[P02-to-P11 affine symbol bridge]
\label{cor:o3af-symbol-bridge}
Under Definition~\ref{def:o3af-ginfty},
\[
 \boxed{
 m_\Gamma(\xi)=c_0+c_1\,q_\Gamma(\xi)
 \qquad\text{with}\qquad
 c_1=1,\quad
 c_0=1-q_\Gamma(0)=1+\log\pi-\psi(\tfrac14).
 }
\tag{O3AF.4}\label{eq:o3af-affine-bridge}
\]
Numerically \(c_0\approx6.372183\), so \(c_0>0\).
\end{corollary}

\begin{proof}
Immediate from~\eqref{eq:o3af-ginfty} and \(m_\Gamma=1+g_\infty\).
\end{proof}

\subsubsection*{Consequences for R31}

\begin{corollary}[R31-B transfers to the concrete P11 operator]
\label{cor:o3af-r31-transfer}
Under Definition~\ref{def:o3af-ginfty} and
Corollary~\ref{cor:o3af-symbol-bridge},
Proposition~\ref{prop:o3ad-gamma-antilocal} applies with \(c_1=1\ne0\).
Hence for every nonzero \(f\in L^2(\mathbb R)\) with
\(\operatorname{ess\,supp}f\subset[-R,R]\), the full-line multiplier action of
\(m_\Gamma\) cannot vanish on any open subinterval of \((R,\infty)\) nor of
\((-\infty,-R)\).

Consequently the annular identity
\eqref{eq:o3ad-forced-cancellation} in R31 holds with the exact
off-diagonal kernel
\[
 K_\Gamma(u)=-\frac{e^{-|u|/2}}{1-e^{-2|u|}},\qquad u\ne0,
\]
and the right-hand side of that identity is real-analytic on each
half-annulus in the concrete P11 setting.  In particular R31-B is now a
theorem about the concrete finite-window operator \(C_{\Gamma,S}\) and no
longer only about the abstract P02 symbol.
\end{corollary}

\begin{proof}
The affine bridge~\eqref{eq:o3af-affine-bridge} places \(m_\Gamma\) exactly in
the range covered by
Proposition~\ref{prop:o3ad-gamma-antilocal}, with \(c_1=1\).  The kernel and
analyticity conclusions are exactly the statements of that proposition and its
proof.
\end{proof}

\begin{remark}[Scope]
\label{rem:o3af-scope}
This subsection closes the manuscript-level interface \(?[O]\) recorded in
Remark~\ref{rem:o3ad-symbol-interface}.  It does not otherwise change the
mathematical status of R30, R30-F, R31, R31-D, R32, or R32-F.  In particular
R30-F remains \(?[O]\).  The R14 firewall is untouched: none of the material
in this subsection makes any polar-gauge, terminal-transport, Object-X, or RH
statement.
\end{remark}
